\documentclass[11pt,a4paper]{article}
\usepackage[utf8]{inputenc}
\usepackage[T1]{fontenc}
\usepackage[french]{babel}
\usepackage{amsmath}
\usepackage{amssymb}
\usepackage{amsthm}
\usepackage{graphicx}
\usepackage{float}
\usepackage{booktabs}
\usepackage{hyperref}
\usepackage{geometry}
\geometry{left=2.5cm,right=2.5cm,top=2.5cm,bottom=2.5cm}

% ============================================================================
% DÉFINITIONS DES ENVIRONNEMENTS
% ============================================================================
\newtheorem{prop}{Proposition}[section]
\newtheorem{lemme}{Lemme}[section]
\newtheorem{thm}{Théorème}[section]
\newtheorem{defi}{Définition}[section]
\newtheorem{cor}{Corollaire}[section]
\newtheorem{rem}{Remarque}[section]
\newtheorem{algo}{Algorithme}[section]
\newtheorem{example}{Exemple}[section]

\newcounter{Indice}
\setcounter{Indice}{0}
\newenvironment{Formule}[1]%
{\begin{eqnarray} \stepcounter{Indice}
		\newcounter{#1}
		\label{#1}
	}%
	{\end{eqnarray}
	\noindent    }

% ============================================================================
% MACROS ET NOTATIONS
% ============================================================================
\def \fcn {{\hat f}_n}
\def \acnm {{\hat \alpha}_n^{(m)}}
\def \acn {{\hat \alpha}_n}
\def \Pcn {{\hat P}_n}
\def \tcn {{\hat \theta}_n}
\def  \ftheta {{f_{\theta}}}
\def \DR {{\cal D}}
\def \AR {{\cal A}}
\def \BR {{\cal B}}
\def \CR {{\cal C}}
\def \LR {{\cal L}}
\def \FR {{\cal F}}
\def \ER {{\cal E}}
\def \HR {{\cal H}}
\def \IR {{\cal I}}
\def \SR {{\cal S}}
\def \Re {{\cal R}}
\def \ftz {{f_{\theta_{0}}}}
\def \fgamma {{f_{\gamma}}}
\def \fetze {{f_{\theta_{0},\varepsilon}}}
\def\RR{{\hbox{${I}$\kern-.22em\hbox{$R$}}}}

\def\proof#1. {\par
	\ifdim\lastskip<15pt
	\removelastskip\penalty-200
	\vskip5pt plus3pt minus3pt
	\fi
	{\def\a{#1}
		\ifx\a\empty
		{\noindent\bf Preuve.}
		\else
		{\noindent\bf Preuve de #1.}
		\fi}\enspace}
\def\endproof{\hfill\hspace{-6pt}\rule[-4pt]{6pt}{6pt}
	\vskip8pt plus3pt minus 3pt}
\def\proofend{\endproof}

\newcommand{\reel}{\mbox{\rm I\hspace{-1mm}R}}
\newcommand{\nat}{\mbox{\rm I\hspace{-1mm}N}}
\newcommand{\un}{\mbox{\rm 1\hspace{-1mm}I}}
\newcommand{\proba}{\mbox{\rm I\hspace{-1mm}P}}
\newcommand{\prob}{$(\Omega,{\cal A},P)$}

% ============================================================================
% DÉBUT DU DOCUMENT
% ============================================================================
\begin{document}
	
	\title{\textbf{Codage Uniforme Fort de Données Temporelles \\ par Processus de Comptage}}
	\author{J. NEMB\'E. \\
		Technical Report N$^\circ$ 003/22/02/2015/MD/JN \\ 
		African Institute of Computing Sciences \\ 
		P.O. Box 2263. Libreville, Gabon.}
	\date{Mars 2026}
	
	\maketitle
	
	\begin{abstract}
		\noindent Les processus temporels tels que la vidéo, la voix ou le son sont encodés en utilisant des modèles stochastiques basés sur des processus de comptage dans le modèle d'intensité multiplicative (ou d'Aalen). Il est démontré que cela produit une technique de codage efficace même si les caractéristiques statistiques du processus sous-jacent ne sont pas prises en compte. Des bornes sont fournies pour la longueur de code la plus courte en moyenne et presque sûrement. Des bornes précises pour la longueur de code et la redondance du code sont dérivées dans le cas où le processus sous-jacent est un processus de Poisson non homogène. Les résultats de simulation montrent que ces bornes sont précises.
		
		\medskip
		\noindent \textbf{Mots-clés :} Codage statistique ; Estimation non paramétrique ; Compression de processus ponctuels ; Modélisation stochastique ; Compression vidéo.
		
		\medskip
		\noindent \textbf{Codes AMS :} 60G55; 62G05; 62G20; 62E25; Secondaire 46E22.
	\end{abstract}
	
	\tableofcontents
	\newpage
	
	% ============================================================================
	\section{Introduction}
	% ============================================================================
	
	Les processus temporels modélisent des flux de données à états finis changeant de manière aléatoire au cours du temps, tels que le son et la vidéo. Par exemple, une vidéo est souvent considérée comme un ensemble d'images (un processus spatial) observé au cours du temps. Dans ce cas, les états sont des couleurs.
	
	Les techniques les plus efficaces pour le codage vidéo et sonore combinent des transformées de type Fourier telles que la DCT, le codage statistique des coefficients (comme Huffman) et parfois la compensation de mouvement comme dans MPEG. Aucune de ces techniques ne considère les trajectoires individuelles des pixels. Au lieu de cela, une compression spatiale est effectuée à intervalles de temps réguliers.
	
	Ce travail présente une technique de compression basée sur l'historique individuel des pixels, mais l'accent n'est pas mis sur la technique de codage elle-même, plutôt sur les performances statistiques de cette technique. D'abord, il est montré que ce problème peut être modélisé par des processus ponctuels dans le modèle d'intensité multiplicative. Cela se fait en considérant les temps de changement de pixel comme des sauts d'un processus ponctuel sous-jacent. L'attribution de codes courts aux temps de saut dans les zones temporelles de plus grande activité devrait produire un code plus court en moyenne.
	
	La procédure d'encodage nécessite des modèles précis pour la distribution des sauts sur l'intervalle de temps $[0,T]$. Dans ce contexte, de nombreux modèles mathématiques basés sur des processus ponctuels peuvent décrire la vidéo sur un intervalle de temps fini. À notre connaissance, aucun modèle de codage vidéo basé sur des processus ponctuels n'a été exposé dans la littérature. Cependant, Brémaud a considéré le codage dynamique de données temporelles avec des processus ponctuels. Pour une revue du modèle d'intensité multiplicative et des propriétés de base des processus ponctuels, voir par exemple Aalen \cite{Aalen75}, \cite{Aalen1978c}, \cite{Aalen1978a}, \cite{Andersen1985}, \cite{Andersen1982a}, \cite{Csorgo1982}, \cite{Campbell1982}, \cite{Anderson1980}.
	
	Puisque les données à encoder consistent principalement en des entiers, de nombreux résultats sur le codage optimal des entiers ainsi que la complexité stochastique qui seront utilisés par la suite peuvent être trouvés dans Rissanen \cite{Rissanen83}, Barron \cite{Barron85}.
	
	% ============================================================================
	\section{Le Modèle Mathématique}
	\label{sec:modele_mathematique}
	% ============================================================================
	
	\subsection{Le modèle d'intensité multiplicative}
	\label{subsec:modele_intensite}
	
	Soit $(\Omega,\FR,P)$ un espace de probabilité et soit ${\FR}_t, t\in [0,T]$ une famille croissante, continue à droite de sous-sigma algèbres de $\FR$. Soit $N$ un processus de comptage tel que $E(N(T)) < \infty$.
	
	Il découle de la décomposition de Doob-Meyer que $N = A + M$ où $A$ est un processus prévisible et $M$ est une martingale. Nous supposerons qu'il existe un processus $\lambda$ non-négatif, continu à gauche, adapté à ${\FR}_t$, avec des limites à droite tel que $A(t) = \int_0^t \lambda(s)ds$.
	
	Alors $M(t) = N(t) - A(t)$ est une martingale de carré intégrable avec processus de variance $<M>(t) = \int_0^t \lambda(s)ds.$
	
	Dans cet article, nous considérons le modèle d'intensité multiplicative (Aalen, 1978) où il est supposé que $\lambda$ est le produit de deux termes :
	\begin{equation}
		\lambda(s) = \alpha(s)Y(s),
	\end{equation}
	où $\alpha(\cdot)$ est une fonction déterministe inconnue appelée fonction d'intensité, tandis que $Y$ est un processus stochastique observable avec $E(Y(s)) = {\theta}_\alpha(s)$.
	
	La fonction $\alpha$ et les trajectoires de $Y$ sont non-négatives et continues à gauche avec des limites à droite, et $Y$ est supposée être adaptée à ${\FR}_t$. Il s'ensuit que :
	\begin{equation}
		dN(s) = \alpha(s)Y(s)ds + dM(s)
	\end{equation}
	et pour tout $t \in [0,T]$ :
	\begin{Formule}{RelN}
		N(t) = \int_0^t \alpha(s)Y(s)ds + M(t),
	\end{Formule}
	et
	\begin{Formule}{EspN}
		E\left [N(t)\right ] = \int_0^t \alpha(s)\theta_{\alpha}(s)ds.
	\end{Formule}
	
	\subsection{Modèles d'encodage temporel}
	\label{subsec:modeles_encodage}
	
	Soit $(\Omega,\FR,P)$ un espace de probabilité et soit ${\FR}_t, t\in [0,T]$ une famille croissante, continue à droite de sous-sigma algèbres de $\FR$. Soit $N$, un processus de comptage tel que $E(N(T)) < \infty$.
	
	Un processus temporel tel que la vidéo ou le son peut être modélisé dans ce cadre comme une séquence de temps discret de processus spatiaux (sources sonores ou écrans vidéo). Dans le cas du codage vidéo, les sources seraient les pixels. Au lieu de coder une vidéo comme une séquence d'images, l'encodage est effectué sur l'histoire de chaque pixel au cours du temps. Une telle histoire peut être modélisée par un processus stochastique à états finis $X(t,\omega)$ modélisant la couleur, ainsi qu'un processus ponctuel $N$, enregistrant les temps de changement.
	
	Après l'observation, un échantillon $((X_1,N_1),\dots,(X_n,N_n))$ est observé sur $[0,T]$. Pour résumer, notre objectif principal est de remplacer le codage spatial usuel basé sur la redondance de fréquence spatiale par un encodage temporel des temps de changement.
	
	Rappelons que l'encodage d'un nombre $x$ nécessite environ $\log K$ bits lorsque $x$ est borné par une constante $K > 1$, et $\log(x) + 2\log\log x$ bits lorsqu'aucune borne n'est supposée (voir \cite{Rissanen83} pour le codage universel d'entiers).
	
	\begin{rem}
		Puisque les temps d'observation sont discrétisés, nous observerons plutôt un processus $N$ qui est obtenu comme une version en temps discret de $\tilde{N}$. Cela signifie qu'entre deux instants $t_1$ et $t_2$, nous ne sommes pas intéressés par le nombre de sauts de $\tilde{N}$, mais seulement si au moins un saut s'est produit, c'est-à-dire que $X_{t_1} \neq X_{t_2}$. Il y a un cas que nous ne traiterons pas, c'est le cas où $X_{t_1} = X_{t_2}$ et que des sauts se sont produits dans l'intervalle de temps. On supposera que c'est un événement rare.
	\end{rem}
	
	\subsection{Modélisation du nombre de sauts avec des processus de comptage}
	\label{subsec:modelisation_sauts}
	
	Ici, on suppose que $n$ sources sont observées sur un intervalle de temps $[0,T]$, échantillonnées à un taux $T/r$ (avec $r$ un entier). Le processus est supposé avoir $m=Card(\SR)$ états (couleurs ou fréquences sonores), de sorte que le modèle sera déterminé par le triplet $(m,n,r)$. Un saut consistera en une paire $(\tau,s) \in [0,r] \times \SR$ d'un index temporel $\tau$ et d'un état $s$.
	
	Afin de simplifier les notations, nous écrirons $N$ pour $N(T)$ ou $N(T,\omega)$. Puisque le temps est discrétisé, à une précision appropriée, le nombre de sauts dans chaque période de temps est une statistique suffisante pour récupérer l'historique global du processus.
	
	L'encodage de base d'un entier borné $m$ nécessite environ $\log m$ bits, et $\log m + 2 \log \log m$ lorsque $m$ n'est pas supposé borné. L'encodage de toute la séquence de $r$ processus spatiaux nécessitera alors :
	\begin{equation}
		\Delta_0=nr\log(m).
	\end{equation}
	
	Cela vient du fait que pour chaque source, l'encodage temporel nécessite $r\log(m)$ bits ($r$ emplacements temporels et $\log m$ bits pour le nouvel état à l'emplacement temporel). Ceci est uniquement un encodage d'état.
	
	Coder uniquement les états aux temps de saut devrait réduire la longueur de code si le nombre de sauts est petit. Soit $q_s$, $q_t$, et $q_c$, des procédures d'encodage d'entiers définies respectivement pour les états de couleur, les temps de saut et le nombre total de sauts.
	
	L'encodage de toutes les données nécessite :
	\begin{equation}
		\sum_{i=1}^n \sum_{j=1}^r q_s(x_{ij}),
	\end{equation}
	où $x_{ij}$ est le $j$-ème saut pour le processus $i$.
	
	Le modèle stochastique le plus général pour ce problème sera de trouver une loi de probabilité $P_{\alpha,\beta}$, où $\alpha$ est une fonction d'intensité inconnue modélisant le comportement temporel et $\beta$ est un paramètre fonctionnel dépendant du temps pour les probabilités d'état.
	
	Soit pour tout $1 \leq i \leq n$, un processus $N_i$ défini par ses comptes totaux de sauts $\mu_i$ et ses emplacements de sauts $\tau_{ij}$. Donc $N_i = (\mu_i ; \tau_{i1},\dots,\tau_{i\mu_i})$. Soit $Y_i = (Y_{i1},\dots,Y_{iN_i{(T)}})$ le processus d'états. Le problème que nous souhaitons modéliser est de déterminer :
	\begin{equation}
		P_{\alpha,f} = \arg \min_{\beta,g} \left [  -\log P_{\beta,g} \left ( (N_1,Y_1),\dots, (N_n,Y_n) \right ) \right ].
	\end{equation}
	
	La première hypothèse que nous ferons est l'indépendance entre le processus temporel et le processus d'état. Si cette hypothèse est vraie, les probabilités pour les temps de saut $Q_t \left ( (N_1,\dots,N_n) \right )$ et les états $Q_s(Y_1,\dots,Y_n)$ peuvent être séparées et optimisées séparément.
	
	L'utilisation de procédures d'encodage statistique telles que Shannon, Fano ou Huffman devrait produire les meilleurs codes en moyenne. En conséquence, l'encodage des données lorsque $(n_i,\tau_i,x_i)$ sont respectivement les résultats observés des comptes de sauts, des emplacements et des états, nécessite :
	\begin{equation}
		q_t((n_1,\tau_1),\dots,(n_n,\tau_n)) + q_s(x_1,\dots,x_n).
	\end{equation}
	
	La tâche principale pour concevoir des codes efficaces est d'étudier des modèles de probabilité alternatifs pour le nombre total de sauts, les échantillons de temps de saut et les échantillons d'états. L'accent dans cet article sera mis sur les modèles pour les échantillons de temps de saut. Plus précisément, comment peut-on modéliser les temps de saut et les comptes associés.
	
	Notre objectif principal n'est pas encore la compression statistique. Même la façon dont un seul processus est encodé peut être optimisée. Ceci est détaillé dans la sous-section suivante.
	
	% ============================================================================
	\section{Encodage Temporel Uniforme}
	\label{sec:encodage_uniforme}
	% ============================================================================
	
	Ici, on supposera que les probabilités pour les états, les comptes de sauts et les emplacements de sauts sont respectivement uniformes sur $[0,m]$, $[0,r]$ et $[0,r]$ puisque le nombre maximum de sauts est $r$, le nombre d'emplacements temporels, et le nombre d'états est $m$.
	
	Dans ce cas, pour tout $1 \leq i \leq n$, $q_s(x_i) = \log m$, $q_t(n_i) = \log r$, $q_t(\tau_i) = \log r$ et la procédure résultante est appelée encodage uniforme de temps et d'état.
	
	L'encodage du nombre de sauts et de leurs emplacements peut être fait de différentes manières.
	
	\begin{enumerate}
		\item \textbf{Vecteur d'états.} Dans le cas où un changement se produit à presque chaque emplacement temporel, il n'y a plus de gain temporel dans la longueur de code. Ainsi, coder le temps devient inutile et seuls les états sont encodés. La longueur de code résultante est :
		\begin{equation}
			\CR_0 = r\log m.
		\end{equation}
		
		\item \textbf{Liste des temps de changement et des états.} Puisque l'encodage d'un changement nécessite sans aucune hypothèse supplémentaire $\log r$ bits pour le temps et $\log m$ bits pour l'état, une façon naturelle d'encoder la séquence est d'encoder le nombre de temps de saut, et pour chaque saut, indiquer le temps et le nouvel état. L'encodage de toute la séquence nécessiterait alors dans le pire des cas :
		\begin{equation}
			\CR_1 = \log(N) + N\log(rm).
		\end{equation}
		
		\item \textbf{Vecteur booléen des temps de changement et liste des états.} Au lieu de coder le nombre de sauts et leurs temps, ce qui a un coût de $N\log r$ bits, on peut encoder un vecteur booléen de $r$ bits avec la signification suivante. Si le bit $i$ est mis à un, cela signifie qu'à la $i$-ème position il y a un saut. Le nombre de uns dans ce champ de bits est aussi le nombre total de sauts. Ainsi, les termes $\log N$ et $N\log r$ dans la longueur de code précédente sont remplacés par $r$, le nombre de bits dans le vecteur booléen des temps de changement. La longueur de code est alors :
		\begin{equation}
			\CR_2 = r + N\log m.
		\end{equation}
		Ceci est commode si le nombre de sauts est grand. Remplacer les temps de sauts encodés en utilisant $\log r$ bits par un seul vecteur booléen de $r$ bits peut produire une longueur de code plus courte.
		
		\item \textbf{Configuration des temps de changement et liste des états.} Au lieu d'effectuer l'encodage des emplacements temporels de manière brute en utilisant $r$ bits comme décrit dans la section précédente, on peut encoder le nombre de sauts $N$, et l'index de la configuration des temps de saut parmi toutes les configurations avec $N$ sauts. Le nombre total de telles configurations est $C_r^N$. Cela donne :
		\begin{equation}
			\CR_3 = \log N + \log C_r^{N} + N\log m.
		\end{equation}
	\end{enumerate}
	
	\begin{prop}
		\label{prop:longueurs_code_uniformes}
		En supposant des probabilités uniformes sur les espaces de temps et d'état, on obtient les longueurs de code suivantes :
		\begin{enumerate}
			\item Encodage d'état complet : $L_1(r,m) = r \log m$
			\item Liste des temps de saut : $L_2(r,m,N) = \log N + N \log(rm)$
			\item Encodage par vecteur booléen : $L_3(r,m,N) = r + N \log m$
			\item Adresse combinatoire : $L_4(r,m,N) = \log N + \log C_r^N + N \log m$,
		\end{enumerate}
		où $N$ représente le nombre de temps de saut (variations d'état).
	\end{prop}
	
	% ============================================================================
	\section{Cadres Statistiques pour l'Estimation d'Intensité}
	\label{sec:cadres_statistiques}
	% ============================================================================
	
	Dans cette section, nous développons quatre cadres statistiques distincts pour la modélisation et l'estimation des processus ponctuels dans le contexte du codage vidéo. Chaque cadre correspond à une structure de données spécifique et nécessite une formulation particulière de la vraisemblance et de l'estimateur de complexité minimale.
	
	\subsection{Cadre 1 : Vecteur de Processus Ponctuels}
	\label{subsec:cadre_vecteur}
	
	\subsubsection{Définition et Modélisation}
	
	Soit $\mathbf{N}(t) = (N_1(t), \dots, N_d(t))^\top$ un vecteur de $d$ processus de comptage indépendants observés sur $[0,T]$. Chaque composante $N_j$ représente un canal d'information (par exemple, les canaux R, G, B d'un pixel, ou plusieurs pixels corrélés).
	
	L'intensité stochastique du vecteur est donnée par :
	\begin{equation}
		\boldsymbol{\lambda}(t) = (\lambda_1(t), \dots, \lambda_d(t))^\top,
	\end{equation}
	où chaque composante suit le modèle multiplicatif d'Aalen :
	\begin{equation}
		\lambda_j(t) = \alpha_j(t)Y_j(t), \quad j = 1, \dots, d.
	\end{equation}
	
	\subsubsection{Hypothèse 1 : Indépendance entre Temps et État}
	
	Sous l'hypothèse d'indépendance entre le processus temporel et les états (couleurs), la vraisemblance se factorise :
	\begin{equation}
		L_{\boldsymbol{\alpha}}(\mathbf{N}) = \prod_{j=1}^d L_{\alpha_j}(N_j),
	\end{equation}
	avec pour chaque composante :
	\begin{equation}
		L_{\alpha_j}(N_j) = \exp \left( \int_0^T \log(\alpha_j(s)) dN_j(s) - \int_0^T \alpha_j(s)Y_j(s) ds \right).
	\end{equation}
	
	\subsubsection{Hypothèse 2 : Lois Uniformes Simplificatrices}
	
	Si on suppose que les intensités sont uniformes sur $[0,T]$ et que les états sont uniformément distribués parmi $m$ couleurs possibles :
	\begin{equation}
		\alpha_j(t) = \frac{E[N_j(T)]}{T \cdot E[Y_j]}, \quad \forall t \in [0,T].
	\end{equation}
	
	La vraisemblance uniforme devient :
	\begin{equation}
		L_{unif}(\mathbf{N}) = \prod_{j=1}^d \left( \frac{1}{T} \right)^{N_j(T)} \exp\left( - \frac{N_j(T)}{T} \cdot T \right) = \prod_{j=1}^d \left( \frac{1}{T} \right)^{N_j(T)} e^{-N_j(T)}.
	\end{equation}
	
	\subsubsection{Estimateur de Complexité Minimale (MDL)}
	
	L'estimateur MDL pour le vecteur de processus est défini par :
	\begin{equation}
		\hat{\boldsymbol{\alpha}}_n = \arg \min_{\boldsymbol{\beta} \in \Gamma_n} \left( -\log L_{\boldsymbol{\beta}}(\mathbf{N}) + C_n(\boldsymbol{\beta}) \right),
	\end{equation}
	où $C_n(\boldsymbol{\beta}) = \sum_{j=1}^d C_n(\beta_j)$ est la pénalité de complexité totale.
	
	Pour une famille de fonctions $\Gamma_n$ de dimension $K$, la pénalité typique est :
	\begin{equation}
		C_n(\boldsymbol{\beta}) = \frac{d \cdot K}{2} \log n.
	\end{equation}
	
	\subsubsection{Espérance de la Longueur de Code}
	
	La longueur de code attendue pour le cadre vectoriel est :
	\begin{equation}
		E[LC_{vect}] = \sum_{j=1}^d \left( E[N_j(T)] \log T + E[N_j(T)] \log m + \frac{K}{2} \log n \right).
	\end{equation}
	
	Sous l'hypothèse uniforme, cela se simplifie en :
	\begin{equation}
		E[LC_{vect}^{unif}] = d \cdot \left( \bar{N} \log(Tm) + \frac{K}{2} \log n \right),
	\end{equation}
	où $\bar{N} = \frac{1}{d}\sum_{j=1}^d E[N_j(T)]$.
	
	\subsection{Cadre 2 : Processus Ponctuel Marqué}
	\label{subsec:cadre_marque}
	
	\subsubsection{Définition et Modélisation}
	
	Un processus ponctuel marqué est défini par une suite de couples $(\tau_k, m_k)_{k \geq 1}$ où $\tau_k$ sont les temps de sauts et $m_k \in \mathcal{E} = \{e_1, \dots, e_m\}$ sont les marques (couleurs).
	
	L'intensité marquée est donnée par :
	\begin{equation}
		\lambda(t, m) = \alpha(t, m)Y(t), \quad m \in \mathcal{E}.
	\end{equation}
	
	\subsubsection{Hypothèse 1 : Indépendance entre Temps et Marque}
	
	Sous l'hypothèse d'indépendance, l'intensité se factorise :
	\begin{equation}
		\lambda(t, m) = \lambda_T(t) \cdot p_M(m),
	\end{equation}
	où $\lambda_T(t)$ est l'intensité temporelle marginale et $p_M(m)$ est la probabilité de la marque $m$.
	
	La vraisemblance devient :
	\begin{equation}
		L_{\alpha, p}(N) = \exp \left( \int_0^T \log(\lambda_T(s)) dN(s) - \int_0^T \lambda_T(s)Y(s) ds \right) \cdot \prod_{k=1}^{N(T)} p_M(m_k).
	\end{equation}
	
	\subsubsection{Hypothèse 2 : Lois Uniformes Simplificatrices}
	
	Si le temps et les marques sont uniformément distribués :
	\begin{equation}
		\lambda_T(t) = \frac{E[N(T)]}{T}, \quad p_M(m) = \frac{1}{m}, \quad \forall m \in \mathcal{E}.
	\end{equation}
	
	La vraisemblance uniforme est alors :
	\begin{equation}
		L_{unif}(N) = \left( \frac{1}{T} \right)^{N(T)} e^{-E[N(T)]} \cdot \left( \frac{1}{m} \right)^{N(T)}.
	\end{equation}
	
	\subsubsection{Estimateur de Complexité Minimale (MDL)}
	
	L'estimateur MDL pour le processus marqué est :
	\begin{equation}
		(\hat{\alpha}_n, \hat{p}_n) = \arg \min_{(\beta, q) \in \Gamma_n \times \mathcal{P}_m} \left( -\log L_{\beta, q}(N) + C_n(\beta) + C_n(q) \right),
	\end{equation}
	où $\mathcal{P}_m$ est le simplexe des probabilités sur $m$ marques.
	
	La pénalité de complexité inclut le terme temporel et le terme des marques :
	\begin{equation}
		C_n(\beta, q) = \frac{K_T}{2} \log n + \frac{m-1}{2} \log n,
	\end{equation}
	où $K_T$ est la dimension du modèle temporel et $(m-1)$ est le nombre de paramètres libres pour la distribution des marques.
	
	\subsubsection{Espérance de la Longueur de Code}
	
	La longueur de code attendue pour le processus marqué est :
	\begin{equation}
		E[LC_{marque}] = E[N(T)] \log T + E[N(T)] \log m + H(M) \cdot E[N(T)] + C_n(\alpha, p),
	\end{equation}
	où $H(M) = -\sum_{i=1}^m p_M(e_i) \log p_M(e_i)$ est l'entropie des marques.
	
	Sous l'hypothèse uniforme ($H(M) = \log m$) :
	\begin{equation}
		E[LC_{marque}^{unif}] = E[N(T)] \log(Tm) + \frac{K_T + m - 1}{2} \log n.
	\end{equation}
	
	\subsection{Cadre 3 : Processus Ponctuel Spatial}
	\label{subsec:cadre_spatial}
	
	\subsubsection{Définition et Modélisation}
	
	Un processus ponctuel spatial est défini sur un domaine $\mathcal{S} \subset \mathbb{R}^2$ (par exemple, une image ou un bloc de pixels). L'intensité spatiale est une fonction $\lambda(s)$, $s \in \mathcal{S}$.
	
	La vraisemblance pour un processus de Poisson spatial est :
	\begin{equation}
		L_{\alpha}(N) = \exp \left( \int_{\mathcal{S}} \log(\alpha(s)) N(ds) - \int_{\mathcal{S}} \alpha(s) ds \right).
	\end{equation}
	
	\subsubsection{Hypothèse 1 : Indépendance Spatiale}
	
	Sous l'hypothèse d'indépendance entre les positions spatiales, le processus se factorise sur une partition $\{A_1, \dots, A_r\}$ de $\mathcal{S}$ :
	\begin{equation}
		L_{\alpha}(N) = \prod_{i=1}^r \exp \left( \log(\alpha_i) N(A_i) - \alpha_i |A_i| \right),
	\end{equation}
	où $\alpha_i$ est l'intensité constante sur la région $A_i$ et $|A_i|$ est l'aire de $A_i$.
	
	\subsubsection{Hypothèse 2 : Loi Uniforme Spatiale}
	
	Si l'intensité est uniforme sur $\mathcal{S}$ :
	\begin{equation}
		\alpha(s) = \frac{E[N(\mathcal{S})]}{|\mathcal{S}|}, \quad \forall s \in \mathcal{S}.
	\end{equation}
	
	La vraisemblance uniforme devient :
	\begin{equation}
		L_{unif}(N) = \left( \frac{1}{|\mathcal{S}|} \right)^{N(\mathcal{S})} \exp\left( - E[N(\mathcal{S})] \right).
	\end{equation}
	
	\subsubsection{Estimateur de Complexité Minimale (MDL)}
	
	L'estimateur MDL pour le processus spatial est :
	\begin{equation}
		\hat{\alpha}_n = \arg \min_{\beta \in \Gamma_n} \left( -\log L_{\beta}(N) + C_n(\beta) \right).
	\end{equation}
	
	Pour une représentation par ondelettes 2D ou histogrammes 2D de dimension $K$ :
	\begin{equation}
		C_n(\beta) = \frac{K}{2} \log n + \lambda \cdot \text{pen}_2(\beta),
	\end{equation}
	où $\text{pen}_2(\beta)$ est une pénalité de régularisation spatiale (ex: norme TV ou $L_2$).
	
	\subsubsection{Espérance de la Longueur de Code}
	
	La longueur de code attendue pour le processus spatial est :
	\begin{equation}
		E[LC_{spatial}] = E[N(\mathcal{S})] \log |\mathcal{S}| + E[N(\mathcal{S})] \log m + C_n(\alpha).
	\end{equation}
	
	Sous l'hypothèse uniforme :
	\begin{equation}
		E[LC_{spatial}^{unif}] = E[N(\mathcal{S})] \log(|\mathcal{S}| \cdot m) + \frac{K}{2} \log n.
	\end{equation}
	
	\subsection{Cadre 4 : Processus Markovien de Comptage de Transitions}
	\label{subsec:cadre_markov}
	
	\subsubsection{Définition et Modélisation}
	
	Soit $X(t)$ un processus de Markov à espace d'états fini $E = \{1, \dots, M\}$. On observe les transitions entre états via les processus de comptage $N^{hj}(t)$ qui comptent les transitions de l'état $h$ vers l'état $j$.
	
	L'intensité de transition est :
	\begin{equation}
		\lambda^{hj}(t) = Y^h(t) \alpha^{hj}(t),
	\end{equation}
	où $Y^h(t) = \mathbb{I}_{\{X(t^-) = h\}}$ est le processus à risque (indicateur d'être dans l'état $h$ avant la transition).
	
	\subsubsection{Hypothèse 1 : Indépendance entre Temps et Transitions}
	
	Sous l'hypothèse d'indépendance, les intensités de transition ne dépendent pas du temps :
	\begin{equation}
		\alpha^{hj}(t) = \alpha^{hj}, \quad \forall t \in [0,T].
	\end{equation}
	
	La vraisemblance pour un individu est (cf. \cite{Aalen75}) :
	\begin{equation}
		L_{\gamma}(T) = \prod_{h \neq j} \exp \left( \int_0^T (1 - \gamma^{hj}(s))Y^h(s)ds + \int_0^T \log \gamma^{hj}(s)dN^{hj}(s) \right).
	\end{equation}
	
	Pour $n$ individus indépendants :
	\begin{equation}
		L_{\gamma}^{(n)}(T) = \prod_{h \neq j} \exp \left( \int_0^T (1 - \gamma^{hj}(s))Y_n^h(s)ds + \int_0^T \log \gamma^{hj}(s)dN_n^{hj}(s) \right),
	\end{equation}
	où $Y_n^h = \sum_{i=1}^n Y^{h,i}$ et $N_n^{hj} = \sum_{i=1}^n N^{hj,i}$.
	
	\subsubsection{Hypothèse 2 : Lois Uniformes Simplificatrices}
	
	Si les transitions sont uniformément distribuées parmi les $M(M-1)$ transitions possibles :
	\begin{equation}
		\alpha^{hj} = \frac{E[N^{hj}(T)]}{\int_0^T Y^h(s) ds}, \quad \forall h \neq j.
	\end{equation}
	
	La vraisemblance uniforme devient :
	\begin{equation}
		L_{unif}^{(n)} = \prod_{h \neq j} \left( \frac{1}{T_{obs}^h} \right)^{N_n^{hj}(T)} \exp\left( - N_n^{hj}(T) \right),
	\end{equation}
	où $T_{obs}^h = \int_0^T Y_n^h(s) ds$ est le temps total d'observation dans l'état $h$.
	
	\subsubsection{Estimateur de Complexité Minimale (MDL)}
	
	L'estimateur MDL pour le processus markovien est :
	\begin{equation}
		\hat{\boldsymbol{\alpha}}_n = \arg \min_{\boldsymbol{\gamma} \in \Gamma_n} \left( -\log L_{\boldsymbol{\gamma}}^{(n)}(T) + C_n(\boldsymbol{\gamma}) \right).
	\end{equation}
	
	La pénalité de complexité pour la matrice de transition $M \times M$ est :
	\begin{equation}
		C_n(\boldsymbol{\gamma}) = \frac{M(M-1) \cdot K_{trans}}{2} \log n,
	\end{equation}
	où $K_{trans}$ est la dimension du modèle pour chaque intensité de transition.
	
	\subsubsection{Espérance de la Longueur de Code}
	
	La longueur de code attendue pour le processus markovien est :
	\begin{equation}
		E[LC_{markov}] = \sum_{h \neq j} \left( E[N^{hj}(T)] \log T_{obs}^h + E[N^{hj}(T)] \log m + C_n(\alpha^{hj}) \right).
	\end{equation}
	
	Sous l'hypothèse uniforme :
	\begin{equation}
		E[LC_{markov}^{unif}] = \bar{N}_{trans} \cdot M(M-1) \cdot \log(T_{obs} \cdot m) + \frac{M(M-1) \cdot K_{trans}}{2} \log n,
	\end{equation}
	où $\bar{N}_{trans}$ est le nombre moyen de transitions par paire d'états.
	
	\subsection{Tableau Synthétique des Cadres Statistiques}
	\label{subsec:tableau_synthese}
	
	\begin{table}[h]
		\centering
		\caption{Comparaison des 4 Cadres Statistiques}
		\label{tab:cadres_stats}
		\begin{tabular}{|l|c|c|c|c|}
			\hline
			\textbf{Cadre} & \textbf{Vecteur} & \textbf{Marqué} & \textbf{Spatial} & \textbf{Markov} \\
			\hline
			Dimension & $d$ canaux & $m$ marques & $|\mathcal{S}|$ pixels & $M$ états \\
			\hline
			Intensité & $\boldsymbol{\lambda}(t)$ & $\lambda(t,m)$ & $\lambda(s)$ & $\lambda^{hj}(t)$ \\
			\hline
			Indépendance & Entre canaux & Temps/Marque & Positions spatiales & Transitions \\
			\hline
			Uniforme & $L_{unif} = \prod T^{-N_j}$ & $L_{unif} = (Tm)^{-N}$ & $L_{unif} = |\mathcal{S}|^{-N}$ & $L_{unif} = \prod (T_{obs}^h)^{-N^{hj}}$ \\
			\hline
			Pénalité MDL & $\frac{dK}{2}\log n$ & $\frac{K_T+m-1}{2}\log n$ & $\frac{K}{2}\log n$ & $\frac{M(M-1)K_{trans}}{2}\log n$ \\
			\hline
			$E[LC]^{unif}$ & $d\bar{N}\log(Tm)$ & $N\log(Tm)$ & $N\log(|\mathcal{S}|m)$ & $\bar{N}_{trans}M^2\log(T_{obs}m)$ \\
			\hline
		\end{tabular}
	\end{table}
	
	\subsection{Critère de Sélection du Cadre}
	\label{subsec:critere_selection}
	
	Le choix du cadre statistique optimal se fait par minimisation de la longueur de code MDL :
	\begin{equation}
		\text{Cadre}^* = \arg \min_{\mathcal{C} \in \{\text{Vect}, \text{Marque}, \text{Spatial}, \text{Markov}\}} \left( -\log L_{\mathcal{C}} + C_n(\mathcal{C}) \right).
	\end{equation}
	
	En pratique, ce choix dépend de :
	\begin{enumerate}
		\item La structure des corrélations dans les données (spatiales, temporelles, entre canaux)
		\item Le nombre d'états/couleurs $m$
		\item La densité de sauts $N/T$
		\item L'homogénéité spatiale $H_s$
		\item La régularité temporelle des transitions
	\end{enumerate}
	
	\begin{rem}
		Le cadre marqué est équivalent au cadre vectoriel monochromatique lorsque l'on décompose le processus en $m$ processus simples $N^{(c)}$, un par couleur. Dans ce cas, le terme $N \log m$ disparaît du codage temporel et est remplacé par un coût fixe de carte de couleurs codé une fois par bloc.
	\end{rem}
	
	\begin{prop}
		Sous les hypothèses de régularité standards (voir \cite{versatile_main.tex}), les estimateurs MDL définis dans chaque cadre satisfont la borne de convergence en distance de Hellinger :
		\begin{equation}
			H^2(P, \hat{P}_n) \leq \min_{Q \in \Gamma_n} \left[ H^2(P, Q) + \frac{1}{n} C_n(Q) \right],
		\end{equation}
		$P$-presque sûrement pour $n$ assez grand.
	\end{prop}
	
	% ============================================================================
	\section{Analyse Ponctuelle}
	\label{sec:analyse_ponctuelle}
	% ============================================================================
	
	Les termes $N$ apparaissant dans les quantités précédentes sont des variables aléatoires comptant le nombre d'événements temporels. Ces nombres peuvent être modélisés comme les résultats de processus de comptage temporels, et peuvent être comparés ponctuellement.
	
	Pour les trois techniques de codage temporel décrites dans la section précédente, on suppose que les changements observés au cours du temps sont les résultats d'un processus ponctuel. La principale différence entre ces techniques réside dans la façon d'encoder ces résultats. Une analyse ponctuelle montre que la troisième procédure de codage produit une longueur de code plus courte que les deux premières dans tous les cas. C'est ce que stipule la proposition \ref{prop00}.
	
	\begin{prop}
		\label{prop00}
		$P_{\alpha} \left ( \CR_3 \leq \min(\CR_1,\CR_2) \right )=1$.
	\end{prop}
	
	\proof.
	Ce résultat sera démontré en deux étapes. D'abord, pour tout $N$, prouvons que $\CR_3 \leq \CR_1$ pour tout $N$. Si $\CR_1 < \CR_3$, alors :
	$$\log N + N\log(rm) < \log N + \log(C_r^N) + N\log m.$$
	Cela signifie que $N\log r < \log(C_r^N)$, ou de manière équivalente que $r^N < C_r^N$. Il est facile de vérifier que cela n'est jamais vrai. Donc $\CR_3 \leq \CR_1$ pour tout $N$.
	
	La deuxième affirmation est que :
	$$\log N + \log C_r^N + N\log m \leq r + N\log m,$$
	ainsi que :
	$$\log N + \log C_r^N \leq r,$$
	ou de manière équivalente que :
	$$NC_r^N \leq 2^r,$$
	et cela est toujours vrai.
	\endproof
	
	En utilisant le résultat précédent, l'encodage temporel sera plus performant que l'encodage spatial si $\CR_3 < \CR_0$. Cela signifie après un calcul simple que :
	\begin{equation}
		C_r^N < \frac{m^{r-N}}{N}.
	\end{equation}
	
	Lorsque le nombre de sauts satisfait cette relation, l'encodage temporel est meilleur. Afin de déterminer l'ensemble des résultats pour lesquels cette relation est satisfaite, il suffit de déterminer :
	\begin{equation}
		\Gamma_{m,r} = \left \{ k \in [1..r] : C_r^k < \frac{m^{r-k}}{k} \right \}.
	\end{equation}
	
	La proposition \ref{MainProp} énonce les conditions sous lesquelles l'encodage temporel est meilleur que l'encodage d'état. Cela nécessite des lemmes techniques qui donnent des bornes pour $\Gamma_{m,r}$.
	
	\begin{prop}
		Considérons un modèle avec $m$ états et $r$ emplacements temporels. Soit $g(k)=kC_r^km^k$, et $\gamma_{m,r} = \frac{r}{1 + 1/m}$.
		\begin{enumerate}
			\item[1)] Si $g(\gamma_{m,r}) < m^r$ alors $\Gamma_{m,r}= ]0..r[$.
			\item[2)] Si $g(\gamma_{m,r}) > m^r$ alors on peut trouver $\gamma_0 < \gamma_{m,r}$ et $\gamma_1 > \gamma_{m,r}$ tels que :
			$\gamma_0 = \max_{k < \gamma_{m,r}} \left \{ k : g(k) < m^r  \right \}$ et $\gamma_1 = \min_{k > \gamma_{m,r}} \left \{ k : g(k) < m^r \right \}$. Donc $\Gamma_{m,r}= ]0..\gamma_0[ \bigcup ]\gamma_1,r[$.
		\end{enumerate}
	\end{prop}
	
	\proof.
	Rappelons que :
	\begin{equation}
		\Gamma_{m,r} = \left \{ k \in [1..r] : kC_r^km^k < m^r \right \}.
	\end{equation}
	
	Soit $g(k)=kC_r^km^k$. Pour $1 \leq k \leq \lfloor r/2 \rfloor$, $g$ est une fonction croissante de $k$, mais afin d'étudier les variations pour $k \in ]r/2,r[$ une analyse plus fine doit être faite. Considérons $\frac{g(k+1)}{g(k)}$.
	
	Après un calcul simple, nous obtenons :
	\begin{eqnarray}
		\frac{g(k+1)}{g(k)} = m \frac{r-k}{k}.
	\end{eqnarray}
	
	Cela montre que $g$ est croissante si $k < \frac{r}{1 + 1/m}$. Soit $\gamma_{m,r} = \frac{r}{1 + 1/m}$, le maximum de $g$. Deux cas sont possibles.
	\begin{enumerate}
		\item[1)] Si $g(\gamma_{m,r}) < m^r$ alors $\Gamma_{m,r}= ]0..r[$.
		\item[2)] Si $g(\gamma_{m,r}) > m^r$ alors on peut trouver $\gamma_0 < \gamma_{m,r}$ et $\gamma_1 > \gamma_{m,r}$ tels que :
		$\gamma_0 = \max_{k < \gamma_{m,r}} \left \{ k : g(k) < m^r  \right \}$ et $\gamma_1 = \min_{k > \gamma_{m,r}} \left \{ k : g(k) < m^r \right \}$. Donc $\Gamma_{m,r}= ]0..\gamma_0[ \bigcup ]\gamma_1,r[$.
	\end{enumerate}
	\endproof
	
	\begin{thm}
		\label{MainProp}
		Considérons un modèle avec $m$ états et $r$ emplacements temporels. Soit $g(k)=kC_r^km^k$, et $\gamma_{m,r} = \frac{r}{1 + 1/m}$. Soit $\gamma_{m,r}^0$ et $\gamma_{m,r}^1$ définis comme :
		$\gamma_{m,r}^0 = \max_{k < \gamma_{m,r}} \left \{ k : g(k) < m^r  \right \}$ et $\gamma_{m,r}^1 = \min_{k > \gamma_{m,r}} \left \{ k : g(k) < m^r \right \}$.
		
		La stratégie de codage la plus efficace est d'effectuer un encodage temporel chaque fois que le nombre de sauts satisfait $n_i \in ]0..\gamma_{m,r}^0] \bigcup [\gamma_{m,r}^1,r]$, et d'effectuer un encodage d'état sinon.
	\end{thm}
	
	Le résultat suivant donne des bornes sur la probabilité que l'encodage temporel ne soit pas choisi, et donne donc des bornes sur la longueur de code totale. Pour un seul processus, cette probabilité est donnée par :
	\begin{equation}
		\epsilon_{m,r} = P \left ( N(T) \in ]\gamma_{m,r}^0,\gamma_{m,r}^1[  \right ).
	\end{equation}
	
	La longueur de code pour un processus $N$ est donnée par :
	\begin{equation}
		LC(m,r) = \epsilon_{m,r} \cdot ( r\log m) + (1 - \epsilon_{m,r} ) \cdot \log \left ( N(T)m^{N(T)}C_r^{N(T)} \right ).
	\end{equation}
	
	% ============================================================================
	\section{Estimation de la Fonction d'Intensité}
	\label{sec:estimation_intensite}
	% ============================================================================
	
	Dans cette section, nous détaillons la procédure d'estimation de la fonction d'intensité $\alpha(t)$ pour différentes familles de fonctions. La méthodologie repose sur la maximisation de la vraisemblance pour une dimension fixée $p$, puis sur la sélection de la dimension optimale via le principe MDL.
	
	\subsection{Formulation Générale du Problème d'Estimation}
	\label{subsec:formulation_generale}
	
	\begin{defi}[Famille de Fonctions d'Intensité]
		\label{def:famille_intensite}
		Soit $\{\phi_j(\cdot)\}_{j=1}^{\infty}$ une base de fonctions sur $[0,T]$. La famille de fonctions d'intensité de dimension $m$ est définie par :
		\begin{equation}
			\Gamma_m = \left\{ \alpha(\cdot) : \alpha(t) = \sum_{j=1}^m \alpha_j \phi_j(t), \quad \alpha_j \in \mathbb{R}^+ \right\},
		\end{equation}
		où les $\alpha_j$ sont les coefficients à estimer.
	\end{defi}
	
	\begin{rem}[Bases Usuelles]
		Les bases $\{\phi_j(\cdot)\}$ peuvent être choisies parmi :
		\begin{itemize}
			\item \textbf{Polynômes :} $\phi_j(t) = t^{j-1}$ (base canonique)
			\item \textbf{Fonctions Trigonométriques :} $\phi_j(t) \in \{\cos(2\pi j t/T), \sin(2\pi j t/T)\}$
			\item \textbf{Histogrammes :} $\phi_j(t) = \mathbb{I}_{[t_{j-1}, t_j]}(t)$ (fonctions indicatrices)
			\item \textbf{Splines :} Fonctions polynomiales par morceaux avec continuité
			\item \textbf{Ondelettes :} Base de Haar, Daubechies, etc.
		\end{itemize}
	\end{rem}
	
	\subsection{Vraisemblance du Processus Ponctuel Réel}
	\label{subsec:vraisemblance_reel}
	
	\begin{prop}[Vraisemblance pour le Cadre Réel]
		\label{prop:vraisemblance_reel}
		Pour un processus ponctuel réel observé sur $[0,T]$ avec intensité $\alpha(t) = \sum_{j=1}^m \alpha_j \phi_j(t)$, la log-vraisemblance est :
		\begin{equation}
			\log L_{\alpha}(N) = \int_0^T \log\left(\sum_{j=1}^m \alpha_j \phi_j(s)\right) dN_n(s) - \int_0^T \left(\sum_{j=1}^m \alpha_j \phi_j(s)\right) Y_n(s) ds.
		\end{equation}
	\end{prop}
	
	\proof.
	Cette formule découle directement du modèle multiplicatif d'Aalen (section \ref{subsec:cadre_vecteur}) avec $\lambda(s) = \alpha(s)Y(s)$. La vraisemblance générale est :
	\begin{equation}
		L_{\alpha}(N) = \exp \left( \int_0^T \log(\alpha(s)) dN(s) - \int_0^T \alpha(s)Y(s) ds \right).
	\end{equation}
	En substituant $\alpha(s) = \sum_{j=1}^m \alpha_j \phi_j(s)$ et en prenant le logarithme, on obtient le résultat.
	\endproof
	
	\subsection{Équations du Maximum de Vraisemblance}
	\label{subsec:equations_mle}
	
	\begin{thm}[Équations du Maximum de Vraisemblance]
		\label{thm:equations_mle}
		Les coefficients $\alpha_1, \dots, \alpha_m$ qui maximisent la vraisemblance satisfont le système non-linéaire :
		\begin{equation}
			\int_a^b \phi_k(s) Y_n(s) ds = \int_a^b \frac{\phi_k(s)}{\sum_{j=1}^m \alpha_j \phi_j(s)} dN_n(s), \quad k = 1, \dots, m.
		\end{equation}
	\end{thm}
	
	\proof.
	On dérive la log-vraisemblance par rapport à chaque composante $\alpha_k$ :
	\begin{equation}
		\frac{\partial \log L_{\alpha}}{\partial \alpha_k} = \int_0^T \frac{\phi_k(s)}{\sum_{j=1}^m \alpha_j \phi_j(s)} dN_n(s) - \int_0^T \phi_k(s) Y_n(s) ds.
	\end{equation}
	En annulant cette dérivée pour tout $k = 1, \dots, m$, on obtient le système.
	\endproof
	
	\begin{cor}[Cas $Y_n(s) = n$ Constant]
		\label{cor:cas_y_constant}
		Lorsque $Y_n(s) = n$ est constant (tous les processus toujours à risque), le système devient :
		\begin{equation}
			n \int_a^b \phi_k(s) ds = \sum_{s \in \text{sauts dans } [a,b]} \frac{\phi_k(s)}{\sum_{j=1}^m \alpha_j \phi_j(s)}, \quad k = 1, \dots, m.
		\end{equation}
	\end{cor}
	
	\subsection{Cas Particulier : Fonction Constante ($p=0$)}
	\label{subsec:cas_fonction_constante}
	
	\begin{prop}[Estimateur pour Intensité Constante]
		\label{prop:estimateur_constant}
		Pour une fonction constante $\alpha(t) = \alpha_0$ (cas $m=1$, $\phi_1(t) = 1$), l'estimateur du maximum de vraisemblance est :
		\begin{equation}
			\hat{\alpha}_0 = \frac{\text{Effectif}([a,b])}{n(b-a)} = \frac{N_n(b) - N_n(a)}{n(b-a)}.
		\end{equation}
	\end{prop}
	
	\proof.
	Pour $m=1$ et $\phi_1(t) = 1$, le système du théorème \ref{thm:equations_mle} devient :
	\begin{equation}
		n \int_a^b 1 \cdot ds = \int_a^b \frac{1}{\alpha_0} dN_n(s).
	\end{equation}
	Ce qui donne :
	\begin{equation}
		n(b-a) = \frac{1}{\alpha_0} (N_n(b) - N_n(a)).
	\end{equation}
	D'où :
	\begin{equation}
		\hat{\alpha}_0 = \frac{N_n(b) - N_n(a)}{n(b-a)} = \frac{\text{Effectif}([a,b])}{n(b-a)}.
	\end{equation}
	\endproof
	
	\begin{rem}[Lien avec le Codage Uniforme]
		L'estimateur constant correspond exactement au cas uniforme des documents \cite{uniform_temp_coding.tex} et \cite{Uniform_temporal.tex}. Les longueurs de code $L_1, L_2, L_3, L_4$ de la proposition \ref{prop:longueurs_code_uniformes} sont calculées avec cet estimateur $\hat{\alpha}_0$.
	\end{rem}
	
	\subsection{Cas Particulier : Fonction Constante par Morceaux (Histogrammes)}
	\label{subsec:cas_histogrammes}
	
	\begin{defi}[Famille des Histogrammes]
		\label{def:histogrammes}
		Soit une partition de $[0,T]$ en $K$ intervalles $I_k = [t_{k-1}, t_k)$, $k=1,\dots,K$. La famille des histogrammes est définie par :
		\begin{equation}
			\Gamma_K^{hist} = \left\{ \alpha(\cdot) : \alpha(t) = \sum_{k=1}^K \alpha_k \mathbb{I}_{I_k}(t) \right\}.
		\end{equation}
	\end{defi}
	
	\begin{prop}[Estimateur pour Histogrammes]
		\label{prop:estimateur_histogramme}
		Pour la famille des histogrammes, l'estimateur du maximum de vraisemblance est donné par :
		\begin{equation}
			\hat{\alpha}_k = \frac{N_n(I_k)}{n |I_k|}, \quad k = 1, \dots, K,
		\end{equation}
		où $N_n(I_k)$ est le nombre de sauts dans l'intervalle $I_k$ et $|I_k| = t_k - t_{k-1}$.
	\end{prop}
	
	\proof.
	Pour les histogrammes, les fonctions de base sont $\phi_k(t) = \mathbb{I}_{I_k}(t)$. Ces fonctions ont la propriété importante :
	\begin{equation}
		\phi_k(t) \phi_j(t) = 0 \quad \text{pour } k \neq j.
	\end{equation}
	Ainsi, pour $t \in I_k$, on a $\sum_{j=1}^K \alpha_j \phi_j(t) = \alpha_k$. Le système du théorème \ref{thm:equations_mle} se découple en $K$ équations indépendantes.
	\endproof
	
	\begin{cor}[Système Linéaire]
		\label{cor:systeme_lineaire}
		Pour la famille des histogrammes, l'estimation se réduit à un système linéaire diagonal :
		\begin{equation}
			\mathbf{A} \boldsymbol{\alpha} = \mathbf{b},
		\end{equation}
		où $\mathbf{A} = \text{diag}(n|I_1|, \dots, n|I_K|)$ et $\mathbf{b} = (N_n(I_1), \dots, N_n(I_K))^\top$.
	\end{cor}
	
	\begin{rem}[Avantage pour le Temps-Réel]
		Le cas des histogrammes est particulièrement adapté à la compression temps-réel car :
		\begin{itemize}
			\item L'estimateur a une forme explicite (pas d'itération nécessaire)
			\item Le calcul est de complexité $O(K)$ où $K$ est le nombre de bins
			\item Les compteurs $N_n(I_k)$ peuvent être mis à jour en temps réel
		\end{itemize}
	\end{rem}
	
	\subsection{Sélection de la Dimension Optimale par MDL}
	\label{subsec:selection_dimension_mdl}
	
	\begin{defi}[Critère MDL pour la Dimension]
		\label{def:critere_mdl_dimension}
		La dimension optimale $m^*$ est sélectionnée par :
		\begin{equation}
			m^* = \arg \min_{m \in \{0, 1, \dots, m_{max}\}} \left( -\log L_{\hat{\alpha}^{(m)}}(N) + C_n(m) \right),
		\end{equation}
		où $C_n(m) = \frac{m}{2} \log n$ est la pénalité de complexité pour $m$ paramètres.
	\end{defi}
	
	\begin{algo}[Séquence Croissante de Dimensions]
		\label{algo:sequence_dimensions}
		\begin{enumerate}
			\item[\textbf{Entrée :}] Données $(N_n, Y_n)$, dimension maximale $m_{max}$
			\item[\textbf{Pour}] $m = 0$ à $m_{max}$ faire :
			\begin{enumerate}
				\item Estimer $\hat{\alpha}^{(m)}$
				\item Calculer la log-vraisemblance $\log L_{\hat{\alpha}^{(m)}}(N)$
				\item Calculer la pénalité $C_n(m) = \frac{m}{2} \log n$
				\item Calculer le score MDL : $S(m) = -\log L_{\hat{\alpha}^{(m)}}(N) + C_n(m)$
			\end{enumerate}
			\item[\textbf{Sortie :}] $m^* = \arg \min_m S(m)$ et $\hat{\alpha}^{(m^*)}$
		\end{enumerate}
	\end{algo}
	
	\begin{thm}[Convergence de l'Estimateur MDL]
		\label{thm:convergence_estimateur}
		Sous les hypothèses de régularité de \cite{versatile_main.tex} (Théorème 1), l'estimateur MDL $\hat{\alpha}_n$ sélectionné par l'algorithme \ref{algo:sequence_dimensions} satisfait :
		\begin{equation}
			H^2(P_{\alpha}, P_{\hat{\alpha}_n}) \leq \min_{m} \left[ H^2(P_{\alpha}, P_{\alpha^{(m)}}) + \frac{1}{n} C_n(m) \right],
		\end{equation}
		$P_{\alpha}$-presque sûrement pour $n$ assez grand.
	\end{thm}
	
	% ============================================================================
	\section{Analyse Statistique}
	\label{sec:analyse_statistique}
	% ============================================================================
	
	\subsection{Bornes d'espérance pour la longueur de code}
	\label{subsec:bornes_esperance}
	
	Rappelons que la longueur de code optimale est :
	\begin{equation}
		\label{optimcode}
		\CR_3 = \sum_{i=1}^n \left[ \log(N_i) + \log C_r^{N_i} + N_i \log m \right ].
	\end{equation}
	
	Le résultat suivant énoncera des bornes en espérance pour $\CR_3$.
	
	\begin{prop}
		\label{prop:bornes_esperance}
		La longueur de code optimale a une espérance et une variance qui satisfont :
		\begin{equation}
			E \left [ \CR_3 \right ] \leq n \log m \cdot E[N] \cdot \log  \left ( E[N] \cdot \left ( E \left ( C_r^N \right ) \right ) \right ),
		\end{equation}
		et
		\begin{equation}
			Var \left [ \CR_3\right ] = n\cdot \left [ (\log m)^2 Var[N] + Var \log \left [NC_r^N \right ] \right ].
		\end{equation}
	\end{prop}
	
	\proof.
	En utilisant $E(\log X) \leq \log E(X)$, on obtient :
	\begin{eqnarray}
		E \left [ \CR_3 \right ] &=& n \sum_{j=1}^{m}  E \left[ \log(N_{1j}) + \log C_r^{N_{1j}}  \right] \\
		&\leq & n \sum_{j=1}^{m} \log  \left ( E \left[ N_{1j} \right ] E \left[ C_r^{N_{1j}} \right] \right ) \\
		&\leq & n \sum_{j=1}^{m} \log  \left ( \left ( r - \sum_{k=1}^r P({\tilde{N}}_1^j (I_k)=0)  \right )\cdot E \left ( C_r^{N_{1j}} \right ) \right ).
	\end{eqnarray}
	\endproof
	
	Rappelons que le premier terme dans la formule, $\log(N_{i})$, est un encodage du nombre de sauts du processus modifié $N_{i}$. Cependant, une longueur de code plus courte peut être produite en encodant soit $N_{ij}$ soit $r - N_{ij}$ selon laquelle de ces deux quantités est plus petite. En ajoutant un bit supplémentaire pour indiquer ce qui est encodé, la longueur de code résultante est alors donnée par :
	\begin{equation}
		\label{optimcodeplus}
		\CR_K = \sum_{i=1}^n  \min \left (\log(N_{i}),\log(r-N_{i}) \right ) + \log C_r^{N_{i}} + N_i\log m.
	\end{equation}
	
	Évidemment $\CR_K \leq \CR_3$. Nous obtenons donc le résultat suivant pour la borne sur la longueur de code :
	
	\begin{prop}
		\label{prop:bornes_optimcodeplus}
		La longueur de code optimale $\CR_K$ a une espérance et une variance qui satisfont :
		\begin{equation}
			E \left [ \CR_K \right ] \leq n \cdot \left( E[\min(\log N, \log(r-N))] + E[\log C_r^N] + E[N]\log m \right),
		\end{equation}
		et
		\begin{equation}
			Var \left [ \CR_K \right ] \leq n \cdot \left( Var[\min(\log N, \log(r-N))] + Var[\log C_r^N] + (\log m)^2 Var[N] \right).
		\end{equation}
	\end{prop}
	
	\subsection{Bornes presque sûres et encodage adaptatif}
	\label{subsec:bornes_presque_sures}
	
	Des bornes presque sûres peuvent être facilement dérivées en utilisant des inégalités de type Hoeffding et la bornitude de $N$. L'inégalité de Bennet sera utilisée comme suit :
	
	\begin{prop}
		\label{prop:bennet}
		Soit $X_1,\dots,X_n$ indépendantes avec $a_k \leq X_k \leq b_k$ pour tout $1 \leq k \leq n$. Soit $S_n=\sum_{i=1}^n X_i$, et $\nu_n=E(S_n)$. Alors pour tout $\epsilon_n > 0$ :
		$$P\left ( \vert S_n - \nu_n \vert \geq \epsilon_n \right ) \leq 2\exp\left(-\frac{2n\epsilon_n^2}{\sum_{k=1}^n(b_k-a_k)^2}\right).$$
	\end{prop}
	
	En supposant que les sources sont indépendantes, des résultats presque sûrs pour la longueur de code moyenne peuvent alors être dérivés.
	
	\begin{prop}
		\label{prop:bornes_presque_sures}
		$P_{\alpha}$-presque sûrement, pour $n$ assez grand, ce qui suit tient pour tout $\delta > 0$ :
		\begin{equation}
			\left \vert \overline{\CR}_3  - \nu \right \vert \leq \frac{(1+\delta)}{2}\frac{\log n}{n}((r-1)\log m + \log C_r^{r/2} + \log r).
		\end{equation}
	\end{prop}
	
	\proof.
	En utilisant l'inégalité de Bennet pour les variables aléatoires $X_k=\log( N_k \cdot C_r^{N_k}) + N_k \log m$, on obtient :
	$$P \left [ \left \vert \overline{\CR}_3  - \nu \right \vert \geq \epsilon_n \right ] \leq 2\exp\left(-\frac{2n\epsilon_n^2}{n(\log r + \log C_r + (r-1)\log m)^2}\right).$$
	
	Si le membre de droite est borné par $n^{-(1+\nu)}$, pour tout $\nu > 0$, alors par le lemme de Borel-Cantelli, cet événement est asymptotiquement impossible. L'affirmation s'ensuit.
	\endproof
	
	\begin{rem}
		Bien que le nombre de couleurs $m$ dans une séquence vidéo soit une valeur prescrite, le choix de la longueur d'échantillon $n$ déterminera le nombre effectif de couleurs. Une caractéristique intéressante serait de construire une procédure qui détermine automatiquement le taux d'échelle optimal $r_n$ étant donné $n$ et par conséquent $m(n)$. Une telle procédure détermine $r_n$ selon $n$ et $m(n)$ de telle manière que la redondance de code moyenne reste bornée lorsque $n$ augmente. Bien sûr, il est difficile de considérer ces résultats comme des asymptotiques, mais pour des flux vidéo très larges tels que la vidéo très haute définition, les longueurs d'échantillon peuvent être suffisamment grandes pour valider la terminologie asymptotique.
	\end{rem}
	
	Le résultat suivant donne le taux d'échelle temporel optimal pour des séquences vidéo de définition croissante.
	
	\begin{prop}
		\label{prop:taux_echelle_optimal}
		Le taux d'échelle temporel optimal pour des séquences vidéo de définition croissante est donné par :
		\begin{equation}
			r_n = o \left ( \frac{n}{m(n)\log n} \right ).
		\end{equation}
	\end{prop}
	
	\proof.
	Rappelons que la borne de redondance est donnée par :
	$$\frac{(1+\delta)}{2}\frac{\log n}{n}((r-1)\log m + \log C_r^{r/2} + \log r).$$
	
	Le terme dominant est d'ordre :
	$$r_n\log m_n\frac{\log n}{n}.$$
	
	Ce terme converge vers zéro si :
	$$r_n/\frac{n}{\log m(n)\log n}$$
	converge vers zéro. D'où l'affirmation.
	\endproof
	
	% ============================================================================
	\section{Exemples de Modèles et Algorithmes}
	\label{sec:exemples_modeles}
	% ============================================================================
	
	Rappelons que puisque $N$ n'est pas le résultat du processus ponctuel continu mais sa version discrétisée dans le temps et afin d'éviter la confusion, le processus continu sera noté $\tilde{N}$ et sa version discrétisée $N$. Donc :
	\begin{equation}
		N(t,\omega)= \sum_{k=1}^r \un_{\{ \tilde{N}(I_k) > 0 \}}(\omega).
	\end{equation}
	
	Ainsi $N(T)$ est une variable aléatoire bornée ($0 \leq N(T) \leq r$) avec moyenne et variance données par :
	\begin{eqnarray}
		E[N] &=& r - \sum_{k=1}^r P(\tilde{N}(I_k)=0), \\
		Var[N] &=& \sum_{k = 0}^r P(\tilde{N}(I_k) = 0) \cdot \left ( 1 + P(\tilde{N}(I_k) = 0) \right ).
	\end{eqnarray}
	
	\subsection{Processus de Poisson}
	\label{subsec:processus_poisson}
	
	Dans cette section, l'hypothèse générale est spécialisée au cas d'un processus de Poisson non homogène. Plus précisément, les données sont générées par des processus de Poisson non homogènes avec intensité sinusoïdale, exponentielle et linéaire. Les résultats sont comparés avec les bornes théoriques.
	
	Dans ce cas, la probabilité est définie pour un ensemble $D$ et une mesure $\nu$ comme :
	\begin{equation}
		P(N(D) = k)=\frac{1}{k!}\exp \left ( -\int_D \alpha(u)du \right ) \cdot \left ( \int_D \alpha(u)du \right )^k.
	\end{equation}
	
	\subsection{Données de durée de vie}
	\label{subsec:durees_vie}
	
	Les modèles de durée de vie avec censure peuvent être traités dans le même cadre en considérant les temps de défaillance comme des sauts d'un processus ponctuel. La vraisemblance avec censure est :
	\begin{equation}
		L_{\lambda} = \prod_{i=1}^n \lambda(X_i)^{\delta_i} \exp\left( -\int_0^{X_i} \lambda(s) ds \right),
	\end{equation}
	où $\delta_i = 1$ si l'événement est observé et $\delta_i = 0$ s'il est censuré.
	
	% ============================================================================
	\section{Simulations Numériques}
	\label{sec:simulations}
	% ============================================================================
	
	\subsection{Longueur de code minimale et espérance empirique}
	\label{subsec:longueur_code_minimale}
	
	Des choix réalistes pour $r$ et $m$ conduisent à des formules théoriques plus simples qui peuvent être comparées aux résultats de simulation numérique.
	
	Le premier ensemble de données compare les longueurs de code pour différents résultats possibles. Il confirme les résultats théoriques sur la longueur de code minimale $L=\min(L_1,L_2,L_3,L_4)$ et sa valeur moyenne.
	
	\begin{table}[h]
		\centering
		\caption{Longueurs de code pour différents paramètres}
		\label{tab:simulations}
		\begin{tabular}{||c|c|c|c||}
			\hline
			$r=256$ & $m = 2$ & $m = 16$ & $m = 256$ \\
			\hline
			$L$ & --- & --- & --- \\
			\hline
			$LC$ & --- & --- & --- \\
			\hline
			$E(L)$ & --- & --- & --- \\
			\hline
			$Var(L)$ & --- & --- & --- \\
			\hline
		\end{tabular}
	\end{table}
	
	% ============================================================================
	\section{Conclusion}
	\label{sec:conclusion}
	% ============================================================================
	
	Ce travail a présenté une technique de compression basée sur la modélisation des processus temporels par des processus ponctuels dans le modèle d'intensité multiplicative. Les principaux résultats sont :
	
	\begin{enumerate}
		\item Quatre cadres statistiques ont été développés pour modéliser différents types de données vidéo (vecteur, marqué, spatial, markovien)
		\item Des bornes théoriques précises ont été dérivées pour les longueurs de code optimales
		\item L'analyse ponctuelle montre que l'encodage combinatoire ($L_4$) est optimal dans la plupart des cas
		\item L'estimation MDL permet de sélectionner automatiquement le modèle optimal
		\item Les simulations confirment que les bornes théoriques sont précises
	\end{enumerate}
	
	Les perspectives incluent l'extension à la compression vidéo haute définition, l'intégration de marques sémantiques, et le développement d'une structure pyramidale pour la navigation multi-échelle.
	
	% ============================================================================
	% BIBLIOGRAPHIE
	% ============================================================================
	\newpage
	\bibliographystyle{alpha}
	\bibliography{biblio}
	
\end{document}